\documentclass[10pt,a4paper]{amsart}
\usepackage[margin=19mm]{geometry}
\usepackage{amsmath,amssymb,mathtools}
\usepackage[scr=rsfs,cal=euler]{mathalfa}
\usepackage{xfrac}
\usepackage[unicode,hidelinks,bookmarks=false]{hyperref}
\newcommand{\ie}{i.e.\ }
\newcommand{\cf}{cf.\ }
\newcommand{\resp}{respectively\ }
\newcommand{\Subsection}{\subsection}
\providecommand{\nobreakdash}{\nobreak\hbox{-}}
\setlength{\emergencystretch}{2em}
\multlinegap0pt
\usepackage{bm}
\usepackage{bbm}
\usepackage{dsfont}
\usepackage{xspace}
\usepackage{cancel}
\usepackage{mathtools}
\usepackage{amsthm}
\makeatletter
\@ifundefined{Theorem}{\newtheorem{Theorem}{Theorem}}{}
\makeatother
\def\End{\mathrm{End}}         
\def\Hom{\mathrm{Hom}}           \def\tenkS{\otimes_{kS}}
\def\soc{\mathrm{soc}}           \def\tenOS{\otimes_{\mathcal{O}S}}
\title[The source permutation module of a block of a finite group a]{The source permutation module of a block of a finite group algebra}
\author{Radha Kessar, Markus Linckelmann}
\date{Source: arXiv:2507.15125v1 (CC BY 4.0)}
\begin{document}
\maketitle

\noindent\emph{Source and licence.} The source permutation module of a block of a finite group algebra. Version arXiv:2507.15125v1 (CC BY 4.0), distributed under the Creative Commons Attribution 4.0 International licence, as recorded in the arXiv licence metadata for this exact version and shown on the abstract page https://arxiv.org/abs/2507.15125v1. Full rights notice: licenses/arxiv-2507.15125-CC-BY-4.0.txt.\par\smallskip

\noindent\textit{Editorial scope.} Counted unit: the Theorem whose source label is \texttt{Brauertree-thm-2} in the article (source0.tex lines 2312--2319, source label \texttt{Brauertree-thm-2}), delivered together with the article's own complete original proof of it (source0.tex lines 2322--2355, one \texttt{proof} environment of 34 lines). No mathematics is written, repaired or re-derived: statement and proof are the article's own text line for line. Required context retained uncounted: none. Every definition, display and label of those lines is retained unchanged. Omitted from this package and not redistributed: 1--2311, 2320--2321, 2356--3414. Nothing inside a retained excerpt is omitted or altered: every retained body is delivered exactly as the article's lines are written, so no omission receipt of any kind accompanies this package. Citations: the retained excerpts carry no citation command at all, so no bibliography entry of the article is transcribed and no reference is redistributed. Reference adaptation: none was needed and none was made; every reference inside the delivered unit resolves inside it, so no unresolved reference or citation is produced. Printed slips kept literal: no printed word, symbol, hypothesis or number of the counted statement or of its proof was altered. Layout and class: the article's own class is \texttt{amsart} with packages including fontenc, amssymb, amsthm, amsmath, xcolor, xy; the wrapper is the lane's pinned \texttt{amsart} editorial wrapper at 10pt on A4, which restates the theorem-like environments the retained text needs and adds the article's own macro definitions that the retained text needs (\texttt{\textbackslash End}, \texttt{\textbackslash Hom}, \texttt{\textbackslash soc}), with extra wrapper packages (bm, bbm, dsfont, xspace, cancel, mathtools). The delivered numbering is the unit's own. Assets: no retained span contains a figure environment, a graphics-inclusion command, a PDF-page inclusion, a table or a TikZ picture; no image file is copied into this package, the package declares no asset dependency and no image file is redistributed with it. Licence: version arXiv:2507.15125v1 of this article is distributed under the Creative Commons Attribution 4.0 International licence, as recorded in the arXiv licence metadata for this exact version and shown on the abstract page https://arxiv.org/abs/2507.15125v1. A full rights notice is delivered at licenses/arxiv-2507.15125-CC-BY-4.0.txt. Characters: every retained excerpt is pure ASCII, so the delivered engine, XeLaTeX with its default Latin Modern text font, reports no missing character.
\begin{Theorem} \label{Brauertree-thm-2}
With the notation above,  the algebra $\End_A(\oplus_{i\in I}\ U_i)$ is a direct 
product of algebras $N_R$, where $R$ runs over the set of
$\langle\rho\rangle$-orbits in $I$, and where $N_R$ is a self-injective Nakayama algebra
such that $\ell(N_R)=$ $|R|$ and such that the projective indecomposable
$N_R$-modules have composition length $|R|$ if $R$ is non-exceptional and
composition length $m|R|$ if $R$ is exceptional.
\end{Theorem} 

\begin{proof}
We keep the notation above. 
If $i$ and $i'$ in $I$ belong to two different $\rho$-orbits, then $U_i$ and $U_{i'}$ have 
no common composition factor, so $\Hom_A(U_i,U_{i'})$ and $\Hom_A(U_{i'}, U_i)$ are
both zero. Thus the endomorphism algebra $\End_A(\oplus_{i\in I} U_i)$ is a direct
product of the algebras $\End_A(\oplus_{s=1}^{a_\rho} U_{\rho^s(i)})$, where
$a_\rho$ is the length of the $\rho$-orbit in $I$, and where $i$ runs over a set
of representatives of the $\rho$-orbits in $I$.

We will show next that there is an $A$-homomorphism $U_{\rho(i)}\to U_i$
with simple kernel and cokernel.
Since $S_{\rho(i)}$ is the top composition factor of $U_i$, there is a surjective
$A$-homomorphism $A_{\rho(i)}\to U_i$. Precomposing this with the
inclusion $U_{\rho(i)}\to A\rho(i)$ yields an $A$-homomorphism
$$\alpha_i : U_{\rho(i)} \to U_i$$
and by comparing composition factors, the image of $\alpha_i$ is the unique
maximal submodule of $U_i$ and the kernel is the unique simple submodule
$\soc(U_{\rho(i)})=$ $\soc(A\rho(i))$ of $U_{\rho(i)}$.

All homomorphisms between the $U_{i'}$ with $i'$ running over the 
$\rho$-orbit of $i$ are linear combinations of compositions of the above.
Thus the quiver of $\End_A(\oplus_{i}\ U_i)$, with $i$ running over a $\rho$-orbit
$R$, is a cycle with $a_\rho=$ $|R|$ edges, and and any full cycle becomes zero if the
$\rho$-orbit is non-exceptional, or $m$-full cycles become zero if the $\rho$-orbit
is exceptional with exceptional multiplicity $m$.
 
It follows that the algebra $\End_A(\oplus_{i}\ U_i)$, with $i$ running over a
$\rho$-orbit $R$, is a split basic  self-injective Nakayama algebra with $a_\rho=$
$|R|$  isomorphism 
classes of simple modules and projective indecomposable modules of 
length $|R| m$ or $|R| $, depending on whether the $\rho$-orbit $R$ of $i$
is exceptional or not.
This proves Theorem \ref{Brauertree-thm-2}.
\end{proof}

\end{document}
